\section{Literature Review}

\subsection{The Macro-Context: From Digital Transformation to Cognitive RPA}

\subsubsection{The Imperative of Digital Transformation in the Modern Economy}
Digital transformation represents the foundational restructuring of enterprise operations, characterized by the deep integration of digital technologies to fundamentally alter how value is delivered, measured, and optimized. In the modern economic paradigm, traditional data silos—where distinct departments hoard isolated datasets and operational metrics—are being rapidly dismantled in favor of interconnected, enterprise-wide digital ecosystems \parencite{peng2022digital}. Historically, artificial intelligence systems deployed within corporate infrastructures have focused almost exclusively on single modalities, such as isolated text-mining algorithms utilized for basic sentiment analysis or standalone image-based recognition models tailored for highly specific, narrowly defined manufacturing applications. However, as the scope of computational capabilities has grown, alongside the exponential increase in big data velocity, so has the critical demand for multimodal data interpretation \parencite{sura2025multimodal}.

The massive accumulation of structured data, including deterministic transaction logs, rigid financial records, and tabular customer relationship management databases, occurs concurrently with the generation of vast lakes of unstructured data. This unstructured matrix encompasses nuanced customer reviews, highly variable visual marketing data, dynamic social media content, and continuous communication audio. The coexistence of these heterogeneous data types necessitates a profound transition from reactive, retrospective analytics to proactive, predictive data-driven intelligence \parencite{gumber2025multimodal}. Empirical analyses indicate that the digital transformation of manufacturing and service enterprises plays a significant, positive role in improving overall enterprise performance, primarily by shifting organizational dependence from traditional, constrained production factors to highly scalable data elements \parencite{peng2022digital}.

By seamlessly integrating these diverse and heterogeneous data formats, enterprises can transcend superficial metrics to achieve a significantly more contextual, robust, and holistic understanding of complex market phenomena \parencite{sura2025multimodal}. Furthermore, evidence suggests that digital transformation indirectly improves performance by systematically optimizing operational efficiency and reducing production costs. Through strategic alignment, enterprises that adopt data-driven innovation ecosystems achieve sustained competitive advantages, fundamentally shifting market power and enabling high-quality growth \parencite{khan2025adaptive}. The ability to process unstructured, multimodal inputs acts as a catalyst for breaking through rigid sectoral barriers, promoting internal specialization, and advancing interdepartmental collaboration, which collectively drive a scalable increase in capital return rates and decision-making efficiency.

\subsubsection{The Evolution of Robotic Process Automation}
Initially, enterprise automation was heavily reliant on traditional Robotic Process Automation (RPA)—software systems explicitly designed to execute highly structured, rule-based, and heavily deterministic tasks. While conventional RPA dramatically improved operational efficiency by automating repetitive, high-volume tasks across standardized user interfaces, its utility was strictly bounded by its inherent fragility \parencite{peng2022digital}. Traditional RPA was fundamentally incapable of processing unstructured data, adapting to even minor user interface modifications, or handling cognitive ambiguities. The prerequisite for successful conventional RPA deployment required highly stable environments with perfectly structured input data, devoid of the variations, noise, and unpredictability typical in human-generated content. Consequently, any deviation from pre-programmed logic resulted in process failure, requiring constant manual intervention and maintenance.

As global supply chains, e-commerce networks, and modern financial systems grow increasingly interconnected and complex, organizations require robust automation frameworks capable of nonlinear pattern recognition and high-dimensional data processing. This imperative has catalyzed the evolutionary leap from standard mechanical automation to "Cognitive RPA," wherein advanced artificial intelligence acts as the central, dynamic reasoning engine \parencite{khan2025adaptive}. In this advanced state, software agents transcend the role of mere mechanical tools to act as quasi-human collaborators on core business tasks, actively participating in data analysis, problem diagnosis, and strategic decision-making \parencite{kephart2021multimodal}.

To be highly effective collaborators in modern corporate environments, these sophisticated systems must move far beyond simple text processing to understand complex multi-modal communications. This involves not only transcribing what is said but deeply interpreting speech complemented by vital non-verbal behaviors and contextual cues, such as deictic gestures, eye gaze, and facial expressions during human-machine interactions \parencite{kephart2021multimodal}. For Cognitive RPA to function efficiently and autonomously, the underlying systems must also satisfy rigorous data quality standards, ensuring the completeness, actuality, accuracy, and validity of the ingested unstructured data.

\begin{table}[htbp]
\centering
\caption{Comparison Between Traditional RPA and Cognitive RPA}
\label{tab:rpa_comparison}
\begin{tabularx}{\textwidth}{@{} l >{\raggedright\arraybackslash}X >{\raggedright\arraybackslash}X @{}}
\toprule
\textbf{Attribute} & \textbf{Traditional RPA} & \textbf{Cognitive RPA} \\
\midrule
Primary Data Input & Highly structured (Relational databases, CSVs) & Unstructured and Multimodal (Audio, Images, Free-text) \\
\addlinespace
Execution Paradigm & Deterministic, rule-based, IF-THEN logic & Probabilistic, learning-based, contextual reasoning \\
\addlinespace
Adaptability & Extremely brittle; fails upon UI or data format changes & Highly adaptive; generalizes across varying inputs and noise \\
\addlinespace
Human Interaction & Minimal; operates strictly in the background & Collaborative; interprets complex cues like gaze and gestures \\
\addlinespace
Tech. Foundation & Screen scraping, basic APIs, macros & Deep Neural Networks, Transformers, Multimodal LLMs \\
\addlinespace
Operational Goal & Labor reduction for repetitive, manual keystrokes & Strategic decision support, anomaly detection, predictive forecasting \\
\bottomrule
\end{tabularx}
\end{table}

\subsubsection{The Landscape of Digital Transformation and RPA in Vietnam}
In the distinct context of Vietnam's rapid digital transformation, local enterprises encounter highly specific operational hurdles. Primarily, they face significant barriers stemming from rigid legacy IT infrastructures and the linguistic complexities of the Vietnamese language, which is tonal, monosyllabic, and characterized by limited large-scale training resources \parencite{le2024phowhisper}. Consequently, attempting to directly apply standard, off-the-shelf Western automation tools often results in systemic failures and unacceptably high error rates during data ingestion \parencite{khan2025adaptive}. This critical gap has catalyzed a pressing demand for localized, multimodal RPA systems. Such bespoke systems increasingly leverage advanced self-supervised learning architectures to effectively process the unique noise inherent in Vietnamese enterprise data—including conversational speech heavily laden with bilingual loanwords and unstructured handwritten documents—while achieving high accuracy at a minimal manual annotation cost \parencite{gumber2025multimodal, zhuo2025vietasr}.

The Fourth Industrial Revolution has exerted a profound impact on the socio-economic landscape of Vietnam, driving digital transformation across critical sectors such as finance, banking, and manufacturing. Within this corporate ecosystem, RPA has emerged as a foundational technology. Unlike previous IT overhauls, such as the implementation of massive Enterprise Resource Planning (ERP) systems—which often demanded exorbitant upfront capital and extensive infrastructural modifications—RPA offers a highly agile and cost-effective alternative. In the Vietnamese context, RPA effectively functions as a "virtual employee," interfacing seamlessly between legacy IT systems and back-office processes to execute repetitive human actions at an accelerated pace and with superior accuracy \parencite{chuong2019rpa}.

A 2019 study underscored that the Vietnamese market presents vast opportunities for RPA deployment, particularly in domains characterized by high transaction volumes and labor-intensive workflows \parencite{chuong2019rpa}. In the banking sector, prominent institutions such as Vietinbank have pioneered the use of RPA to navigate stringent operational processes. For instance, card activation teams utilize RPA bots to validate client requests against complex compliance guidelines, thereby significantly reducing human effort and preventing delays that historically led to customer dissatisfaction. Furthermore, RPA assists Vietnamese commercial banks in account origination by autonomously verifying customer identities, credit scores, and address details. Anti-money laundering (AML) efforts similarly benefit from RPA integration; software agents continuously analyze high-volume transactions against deterministic validation rules, immediately flagging suspicious activities to compliance departments. Beyond banking, RPA is extensively applied across the broader financial services sector in Vietnam to minimize human error and lower operating expenditures. In accounts payable and receivable, the technology frequently employs Optical Character Recognition (OCR) to digitize and extract data from heterogeneous vendor invoices, autonomously routing them for payment processing upon validation \parencite{peng2022digital}.

The manufacturing industry in Vietnam has also adopted RPA to enhance back-office productivity. Production managers deploy RPA agents to monitor operational status by extracting data from Manufacturing Execution Systems (MES) and cross-referencing it with established production standards. A highly critical application within this sector is the automated generation of Bills of Material (BOM). Because errors in a BOM can precipitate inaccurate product costing and severe shipment delays, RPA is utilized to replicate human data-entry steps with high agility and deterministic reliability. A significant case study is Samsung Electronics Vietnam (SEV), which implemented RPA across various departments starting in 2018. Facing the logistical challenge of managing a massive workforce of over 120,000 employees, the company utilized RPA to "virtualize" its back-office operations. This deployment spanned production monitoring—where robots extract hourly data from MES and ERPs to distribute automated reports—as well as purchasing, human resource management, and accounting \parencite{chuong2019rpa}.

Despite these operational triumphs, research indicates key limitations to traditional RPA implementation in Vietnam, emphasizing that deterministic automation is not a universal panacea for all business complexities. While conventional RPA excels at mimicking keystrokes, its internal architecture relies on fixed logic requirements. The technology is strictly restricted to tasks with low cognitive demands that do not necessitate subjective judgment, creativity, or complex interpretation \parencite{kephart2021multimodal}. Because traditional RPA bots operate based on pre-defined scripts, they lack the capacity to adapt to unprogrammed scenarios. In the Vietnamese corporate landscape, where decision-making processes frequently rely on tacit knowledge or non-standardized intuition, this rigidity poses a formidable barrier. Consequently, organizations face the challenge of limited exception handling. When a bot encounters an anomaly—such as a non-standard invoice format or an unexpected software pop-up—it typically fails, necessitating costly manual intervention. Furthermore, localized language identification remains a persistent struggle. The ability of legacy OCR systems to accurately ingest and process the complex Vietnamese language is currently insufficient, especially when processing localized contracts, handwritten notes, or regionally distinct administrative documents \parencite{le2024phowhisper, sura2025multimodal}.

\subsection{Theoretical Concepts: Multimodal AI and Mathematical Foundations}

\subsubsection{Multimodal AI and Cross-Modal Dynamics}
Multimodal Artificial Intelligence refers to computational systems mathematically engineered to ingest, process, align, and reason across heterogeneous data types—or modalities—simultaneously. Drawing biological inspiration from how the human brain synthesizes visual, auditory, and sensory inputs to navigate the physical world, multimodal AI aims to replicate this cohesive perception in silicon. Unlike unimodal systems that are rigidly confined to a single, isolated data format and thus inherently blind to broader contexts, multimodal AI captures richer, multi-dimensional patterns and uncovers latent relationships across multiple, parallel information channels \parencite{khan2025adaptive, sura2025multimodal}.

By fusing textual narratives, visual scenes, auditory frequencies, and numerical time-series inputs into a shared representational space, these models synthesize a comprehensive understanding of complex phenomena. Multimodal systems explicitly exploit the complementary strengths of different data types. For instance, deterministic variables extracted from a numerical database can be contextually enriched and validated by semantic visual and linguistic data processed concurrently, dramatically reducing uncertainty and the false-positive rates inherent to single-channel analysis \parencite{gumber2025multimodal}.

\subsubsection{Acoustic and Speech Processing: Signal Abstraction and Spectral Features}
The processing of continuous auditory data requires a highly precise, lossless transformation of physical acoustic waves into discrete numerical representations suitable for ingestion by deep neural networks. Sound inherently exists as continuous analog waveforms generated by the physical oscillation of air particles over time. To process this mechanical energy via computer systems, the analog signal undergoes Sampling, governed by the Nyquist-Shannon sampling theorem, where the amplitude of the continuous wave is measured at discrete time intervals. Typically executed at frequencies of 16 kHz or 44.1 kHz, this sampling rate dictates the absolute temporal resolution of the captured signal. Sampling is immediately followed by Quantization, an approximation process which maps each sampled amplitude to a finite set of discrete integer values, enabling efficient computational storage and processing.

Raw digital waveforms, while informationally complete, are computationally heavy and contain extremely high acoustic redundancies that offer zero phonetic or semantic value. To extract meaningful linguistic features, speech signals are transformed from the time domain to the frequency domain using the Fast Fourier Transform (FFT). Traditionally, this resulting power spectrum was further processed into Mel-Frequency Cepstral Coefficients (MFCCs) by applying a Discrete Cosine Transform (DCT) to heavily compress and decorrelate the features. However, in modern deep learning paradigms, the DCT step is predominantly omitted to produce Filter Bank features (Fbank or log-Mel spectrograms). Fbank retains the highly correlated spectral structures of speech directly, mirroring the non-linear frequency perception of the human cochlea.

\subsubsection{Visual Processing: Spatial Hierarchies and Attention Maps}
Visual inputs, ranging from scanned enterprise financial documents and automated quality control images to complex product catalog photographs, are initially ingested as massive, unrolled matrices of raw pixel intensities. Through the application of Convolutional Neural Networks (CNNs) or modern Vision Transformers (ViTs), these rigid pixel matrices are mathematically transformed into high-level, abstract visual feature vectors \parencite{khan2025adaptive}.

CNNs utilize spatial convolution operations, applying dynamically learned mathematical filters to detect low-level edges, textures, and eventually hierarchical, complex objects. The fundamental convolution operation is mathematically expressed as:

\begin{equation*}
    y_{ij} = (W * x)_{ij} = \sum_{m} \sum_{n} W_{mn} \cdot x_{(i+m)(j+n)}
\end{equation*}

where  represents the output feature map,  denotes the learnable filter weights (the kernel), and  is the input pixel matrix. To address the vanishing gradient problem in deep networks, advanced CNN architectures utilize residual learning, allowing layers to fit a residual mapping represented by:

\begin{equation*}
    y = F(x, \{W_i\}) + x
\end{equation*}

where  denotes the residual function learned by the stacked layers and  is the identity mapping, ensuring gradients flow directly through shortcut connections.

Conversely, Vision Transformers approach the image by abandoning strict translational invariance. An input image $I \in \mathbb{R}^{H \times W \times C}$ is divided into a grid of non-overlapping sequential patches $x_p \in \mathbb{R}^{N \times (P^2 \times C)}$ and linearly projected into a fixed-dimensional embedding space. Complex self-attention mechanisms probabilistically model the intricate, long-range relationships between different, non-adjacent regions of an image, allowing for highly contextualized visual comprehension.

\subsubsection{Textual and Structured Data Processing}
Structured enterprise inputs, including relational databases, Internet of Things (IoT) sensor outputs, and financial time-series logs, are typically encoded using Multilayer Perceptrons (MLPs) or Recurrent Neural Networks (RNNs). These mathematical models capture complex, non-linear correlations and temporal dependencies embedded deep within numerical configurations, allowing the system to recognize invisible operational trends and forecast long-term systemic trajectories \parencite{gumber2025multimodal}.

Unstructured text processing relies fundamentally on transformer-based architectures, which have revolutionized the extraction of semantic data. Textual strings are tokenized and mathematically mapped into continuous dense vectors residing in a high-dimensional semantic space. The core computational engine of the transformer is the self-attention mechanism, which transforms input embeddings into Query ($Q$), Key ($K$), and Value ($V$) matrices. The attention scores are computed using scaled dot-product similarity:

\begin{equation*}
    Attention(Q, K, V) = softmax \left( \frac{QK^T}{\sqrt{d_k}} \right) V
\end{equation*}

where $d_k$ is the dimension of the keys used to scale the dot product. This calculation evaluates the relative context of every single word against all other words in a given sequence simultaneously, effectively capturing deep linguistic nuance, subtle sentiment, and overarching semantic intent \parencite{dubey2024llama}.

\subsubsection{The LLM-as-Processor Paradigm for Multimodal Alignment}
Traditionally, synthesizing multi-dimensional data required complex cross-modal fusion architectures (e.g., early, late, or joint embedding fusion). However, these mathematical alignment strategies demand immense computational overhead and continuous retraining pipelines, rendering them economically impractical for agile enterprise deployments. 

To overcome this bottleneck, modern RPA architectures are shifting towards the "LLM-as-processor" paradigm. Instead of mathematically concatenating isolated vectors, systems leverage Multimodal Large Language Models (such as the Gemini or LLaMA families) to perform implicit fusion. By ingesting raw audio, images, and text simultaneously through zero-shot prompting and in-context learning, these models bypass the need for explicitly learned fusion layers \parencite{dubey2024llama}. This approach delegates the cognitive alignment of heterogeneous modalities entirely to the LLM's vast pre-trained semantic space, drastically reducing the architectural complexity and operational cost of enterprise RPA nodes.

\subsection{State-of-the-Art Review: Multimodal Encoders, LLMs, and Retrieval Architectures}
This section reviews the theoretical foundations and pioneering studies that form the technological basis for developing a multimodal invoice processing system. The discussion focuses on three core research areas: Vietnamese Automatic Speech Recognition (ASR), information extraction using Large Language Models (LLMs), and hybrid search architectures.

\subsubsection{Vietnamese Automatic Speech Recognition and Feature Representation}
The task of converting Vietnamese speech into text presents several fundamental challenges due to the tonal nature of the language and the diversity of regional dialects. During the preprocessing stage, traditional ASR systems typically rely on Mel-frequency cepstral coefficients (MFCC). However, Filter Bank (Fbank) features have recently become a preferred representation because they preserve richer spectral structures, allowing deep neural networks such as convolutional neural networks and transformer-based architectures to learn phonetic representations directly.

From an architectural perspective, the evolution from Transformer models to Conformer architectures has addressed the challenge of capturing both global contextual dependencies and local phonetic patterns. Conformer integrates self-attention mechanisms with convolutional modules, enabling more effective modeling of speech sequences. More recently, the Zipformer architecture has introduced further improvements by processing speech signals at multiple temporal resolutions, thereby significantly reducing memory consumption and computational cost \parencite{yao2024zipformer}.

In the Vietnamese context, early approaches based on self-supervised learning architectures have marked an important milestone. Models such as Wav2Vec 2.0 demonstrated that high-quality speech recognition can be achieved with significantly less labeled data by leveraging large-scale unlabeled speech corpora \parencite{baevski2020wav2vec}. Building upon the strong foundation of the Whisper model developed by OpenAI \parencite{radford2023robust}, the PhoWhisper architecture introduced by VinAI further improved Vietnamese ASR performance by fine-tuning Whisper on more than 800 hours of multi-regional Vietnamese speech data \parencite{le2024phowhisper}. Experimental results show that PhoWhisper significantly outperforms previous Wav2Vec-based systems in terms of word error rate (WER).

Meanwhile, the open-source VietASR system demonstrates that combining the HuBERT self-supervised learning model \parencite{hsu2021hubert} with a lightweight Zipformer encoder can achieve state-of-the-art performance for low-resource languages while maintaining real-time inference capability \parencite{zhuo2025vietasr}. Therefore, this study adopts the PhoWhisper-base configuration as the core preprocessing module to handle spoken retail transactions containing loanwords and abbreviations.

\textit{Critical Limitation}: Despite achieving state-of-the-art WER for conversational Vietnamese, PhoWhisper inherits the heavy autoregressive decoding bottleneck of the Whisper family; this intrinsic sequential generation mechanism results in high computational latency and massive VRAM requirements that strictly prohibit the sub-second, real-time transcription requirements of high-frequency commercial RPA pipelines. Conversely, while Zipformer and VietASR effectively optimize parameter counts, their reliance on complex multi-stack downsampling and clustering-based SSL pre-computation still incurs non-trivial initialization delays, creating micro-stutters that impede the strictly synchronous streaming requirements mandated by enterprise-grade risk management.

\subsubsection{Zero-shot Information Extraction via Large Language Models}
Traditional Named Entity Recognition (NER) approaches typically rely heavily on manually annotated datasets and lack flexibility when document layouts vary significantly. However, the rapid development of Large Language Models (LLMs), including the Gemini family developed by Google, has introduced a new paradigm for information extraction through zero-shot and few-shot prompting techniques.

Furthermore, advanced multimodal models such as Gemini enable zero-shot information extraction directly from highly unstructured and noisy inputs, including handwritten invoices, scanned documents, or transcribed speech data. In the proposed pipeline, multimodal models are responsible for extracting initial raw entities from heterogeneous data sources. Rather than relying on isolated extraction steps, this architecture utilizes Gemini as the central semantic routing module. Gemini maps raw, unstructured product names from heterogeneous inputs into standardized enterprise database categories without requiring large-scale retraining. Instead of retraining task-specific neural networks, the system defines target entities and domain-specific constraints directly through in-context learning. According to the technical report by \textcite{reid2024gemini}, Gemini's native multimodal architecture demonstrates exceptional instruction-following capabilities and zero-shot reasoning across millions of context tokens, enabling highly effective structured information extraction directly from raw inputs.

\textit{Critical Limitation}: However, deploying end-to-end multimodal LLMs like Gemini to process the entirety of an automation pipeline introduces severe computational bottlenecks; the exorbitant token-generation latency and immense memory footprints render them economically impractical for processing millions of high-frequency micro-transactions typical of continuous RPA workflows.

\subsubsection{Hybrid Search Systems and Vector Databases}
One of the primary limitations of traditional keyword-based retrieval systems arises from the semantic mismatch problem, where synonyms, abbreviations, or spelling variations fail to match exact query terms. Classical search engines such as Elasticsearch rely on sparse retrieval techniques such as TF-IDF or BM25, which often struggle to capture deeper semantic relationships.

In order to address this limitation, dense retrieval approaches have been proposed using transformer-based embedding models such as Sentence-BERT \parencite{reimers2019sentence}. These models map textual inputs into dense vector representations in high-dimensional space, allowing semantic similarity to be measured using metrics such as cosine similarity. However, purely dense retrieval methods may struggle to distinguish between highly similar technical identifiers, for example product specifications such as "LED 8W" and "LED 9W".

To overcome this challenge, hybrid search architectures that combine both sparse and dense retrieval signals have been proposed. Modern vector databases such as Qdrant \parencite{qdrant} support this approach by executing both retrieval strategies in parallel.

\textbf{a. Reciprocal Rank Fusion }

To combine ranking signals from multiple retrieval systems, this study employs the Reciprocal Rank Fusion (RRF) algorithm proposed by \textcite{cormack2009reciprocal}. The RRF scoring function is defined as:

\begin{equation*}
    RRFScore(d) = \sum_{r \in R} \frac{1}{k + r(d)}
\end{equation*}

where $d$ represents a retrieved document, $R$ denotes the set of ranking systems, $r(d)$ represents the rank of document $d$ in each system, and $k$ is a smoothing constant.

\textbf{b. Weighted Score Fusion }

In addition to RRF, this study implements a Weighted Score Fusion strategy. In this approach, the raw relevance scores from both dense and sparse retrieval pipelines are first independently normalized using min-max normalization to ensure comparable scales. The final fused score for each candidate document is then computed as a linear combination:

\begin{equation*}
    Score_{fused}(d) = \alpha \cdot \hat{s}_{dense}(d) + (1 - \alpha) \cdot \hat{s}_{sparse}(d)
\end{equation*}

where $\hat{s}_{dense}(d)$ and $\hat{s}_{sparse}(d)$ denote the min-max normalized scores from the dense and sparse retrieval systems respectively, and $\alpha \in [0, 1]$ is a tunable weight parameter that controls the relative contribution of semantic similarity versus lexical matching. A higher $\alpha$ prioritizes semantic understanding, while a lower $\alpha$ emphasizes exact keyword matching—a critical consideration for domains with dense technical product codes.

\textbf{c. Weighted Reciprocal Rank Fusion}

To further enhance flexibility, this study also proposes a Weighted Reciprocal Rank Fusion (Weighted RRF) strategy that combines the rank-based scoring mechanism of standard RRF with tunable importance weights for each retrieval pipeline. Unlike standard RRF, which assigns equal contribution to all ranking systems, Weighted RRF allows the system to explicitly control the relative influence of dense and sparse retrieval signals. The fused score is defined as:

\begin{equation*}
    Score_{wrrf}(d) = \alpha \cdot \frac{1}{k + r_{dense}(d)} + (1 - \alpha) \cdot \frac{1}{k + r_{sparse}(d)}
\end{equation*}

where $r_{dense}(d)$ and $r_{sparse}(d)$ denote the rank of document $d$ in the dense and sparse retrieval systems respectively, $k$ is a smoothing constant, and $\alpha \in [0, 1]$ is a tunable weight parameter. This formulation inherits the robustness of rank-based fusion from RRF—eliminating the need for explicit score normalization—while simultaneously providing the configurability of weighted fusion. When $\alpha = 0.5$, Weighted RRF reduces to standard RRF with equal contributions.

\textit{Critical Limitation:} Although Weighted RRF elegantly balances lexical and semantic ranking without strict normalization, it inevitably requires multiple, parallel search pipelines to conclude completely before the aggregation mathematics can commence, engendering re-ranking latency that actively breaches the deterministic, sub-second execution thresholds required by real-time enterprise RPA workflows.

\begin{table}[H]
\centering
\small
\caption{Benchmark of State-of-the-Art Models and Enterprise RPA Deployment Constraints}
\label{tab:sota_rpa_gemini}
\begin{tabularx}{\textwidth}{@{} l l >{\raggedright\arraybackslash}X >{\raggedright\arraybackslash}X @{}}
\toprule
\textbf{Modality/Domain} & \textbf{SOTA Model} & \textbf{Core Architectural Innovation} & \textbf{Enterprise RPA Limitation} \\
\midrule
Acoustic (ASR) & PhoWhisper & Supervised weak-scale pre-training fine-tuned on Vietnamese data & High autoregressive decoding latency; massive VRAM consumption prevents edge deployment. \\
\addlinespace
Acoustic (ASR) & VietASR & U-Net-like downsampling, BiasNorm, HuBERT SSL pre-training & Initialization delays and multi-stack complexity cause micro-stutters during real-time streaming. \\
\addlinespace
Extraction (LLM) & Gemini Family & Native multimodal architecture, direct image/audio/text processing, zero-shot semantic routing & Cloud API dependency (security risks), exorbitant token-generation cost and latency impractical for high-frequency closed-loop transactions. \\
\addlinespace
Retrieval (Dense) & Sentence-BERT & Siamese BERT networks for continuous dense vector space representations & Fails to distinguish minute numerical variations in dense technical product codes. \\
\addlinespace
Retrieval (Hybrid) & Qdrant + WRRF & Parallel sparse/dense querying with tunable Weighted Reciprocal Rank Fusion & Pipeline synchronization during re-ranking introduces latency that breaches strict deterministic SLAs. \\
\bottomrule
\end{tabularx}
\end{table}

\subsection{Critical Analysis and The Literature Gap}
\subsubsection{The Fundamental Trade-off: Accuracy, Latency, and Operational Cost}
Although Robotic Process Automation has proven highly effective in improving operational efficiency, the reality of global retail operations dictates that input data is predominantly unstructured and heavily noisy \parencite{gumber2025multimodal}. The inability of conventional RPA systems to cognitively handle such data creates a massive bottleneck for enterprise-scale digital transformation, forcing organizations to rely on costly manual intervention for exception handling \parencite{khan2025adaptive}.

To address this profound limitation, the integration of Artificial Intelligence into automation pipelines has been extensively researched. Nevertheless, a critical synthesis of the current state-of-the-art literature reveals a fundamental, unresolved trilemma between prediction accuracy, computational latency, and operational cost. Traditional sparse retrieval systems provide extremely fast response times and near-zero operational costs, but their accuracy deteriorates precipitously when handling semantically ambiguous data \parencite{reimers2019sentence}.

Conversely, end-to-end multimodal Large Language Models (like the Gemini family) offer superior accuracy and unprecedented zero-shot adaptability. Yet, deploying these massive cloud-based architectures introduces severe computational latency and exorbitant API costs, rendering them entirely impractical for high-frequency RPA ecosystems \parencite{reid2024gemini}.


\subsubsection{Addressing the Gap: An Input-Level Multimodal Pipeline for the Vietnamese Enterprise Context}
A critical analysis indicates a dichotomy in current methodologies: they optimize either for computational efficiency at the expense of semantic accuracy, or for deep multimodal understanding at the cost of operational latency. Compounding this architectural trade-off is the linguistic barrier. Standard, off-the-shelf Western automation tools—typically pre-trained on clean, high-resource languages like English—fail to generalize when applied to Vietnam's localized enterprise environment. Consequently, attempting to directly apply these systems to Vietnamese commercial data results in unacceptably high error rates during ingestion \parencite{khan2025adaptive}. For Vietnamese commercial banks, large-scale retail chains, and national telecommunications firms to fully and reliably automate customer interactions and internal operational workflows, they require bespoke RPA systems capable of accurately processing localized, heavily noisy, and highly unstructured multimodal data \parencite{peng2022digital}. This includes successfully parsing conversational speech containing extensive bilingual loanwords, or digitizing and comprehending handwritten financial documents \parencite{gumber2025multimodal}.

This critical market necessity drives the ongoing strategic shift towards advanced self-supervised learning architectures and highly efficient encoding pipelines that can achieve enterprise-grade accuracy with only minimal manually labeled data \parencite{zhuo2025vietasr}. Meeting these rigorous technical constraints sets the foundational, non-negotiable requirement for the deployment of modern intelligent enterprise solutions within the Vietnamese economic landscape \parencite{sura2025multimodal}.

To explicitly address this significant Literature Gap, future research and system architectures must propose a decoupled, input-level multimodal pipeline optimized for the enterprise edge. By leveraging highly efficient, specialized acoustic encoders like PhoWhisper \parencite{le2024phowhisper} or VietASR \parencite{zhuo2025vietasr} strictly for the initial ingestion of conversational data, and relegating multimodal LLMs like Gemini \parencite{reid2024gemini} exclusively to flexible, zero-shot entity extraction rather than end-to-end processing, systems can aggressively preserve computational resources.. Furthermore, by employing specialized semantic resolution modules—such as hierarchical LLM classifiers and Qdrant-backed \parencite{qdrant} hybrid vector search utilizing Weighted Score Fusion or WRRF \parencite{cormack2009reciprocal}—as the alignment core, enterprises can effectively optimize computational latency and operational cost while maintaining state-of-the-art accuracy for noisy, low-resource multimodal data.





